% ED26-P17-judgment-aggregation-0001. Exact source: List--Pettit,
% author preprint revised 9 October 2001, not an inferred journal version.
\chapter{Judgment Aggregation and Collective Consistency}
\label{sec:p17-judgment-aggregation}

A group may accept each of two premises by a majority and reject their
conjunction by another majority. Separate votes can settle each question even
though the answers do not jointly describe a possible situation. Judgment
aggregation studies which requirements on the inputs,
procedure and collective answers can be maintained together.

The presentation here follows the complete author preprint of List and Pettit,
revised 9 October 2001 \cite{ListPettit2001JudgmentPreprint}. Its appendix gives
Christian List's impossibility proof. The arguments below retain that proof's
separate even- and odd-electorate cases and also prove the source's positive
claims about restricted profiles, priority and special support. A judgment is
binary assent to a proposition. It is neither a degree of confidence nor an
ordering of alternatives. Preference aggregation has the different input and
output types defined in Chapter~\ref{sec:p17-social-choice}.

\section{Agendas, assent and logical closure}
\label{sec:p17-ja-agendas}

\begin{hbdefinition}{def:p17-ja-agenda}{An agenda with opposing propositions}
Let \(N=\{1,\ldots,n\}\), where \(n\geq2\), be a finite electorate. Let
\(X\) be a set of formulas of classical propositional logic. An \emph{agenda}
in the source's sense satisfies the following opposing-proposition condition:
for every \(\varphi\in X\), either \(\neg\varphi\in X\), or
\(\varphi=\neg\psi\) for some \(\psi\in X\).
For the impossibility theorem, \(X\) contains distinct atoms \(p,q\) and
\(p\land q\). Consequently it also contains \(\neg p,\neg q\) and
\(\neg(p\land q)\).

The agenda need not be finite and need not contain every propositional formula.
In particular, the opposing-proposition condition does not require indefinitely
adding iterated negations. The six-formula agenda
\[
 X_0=\{p,\neg p,q,\neg q,p\land q,\neg(p\land q)\}
\]
already satisfies the assumptions.
\index{judgment aggregation!agenda}\index{agenda!logical dependence}
\end{hbdefinition}

\begin{hbdefinition}{def:p17-ja-rational}{Three distinct requirements on a judgment set}
A \emph{judgment set} is a subset \(J\subseteq X\), listing precisely the
propositions assented to. It is:
\begin{enumerate}
\item \emph{complete} if, for every \(\varphi\in X\), at least one of
      \(\varphi\in J\), \(\neg\varphi\in J\), or
      \(\psi\in J\) with \(\varphi=\neg\psi\) holds;
\item \emph{weakly consistent} if no formula and its literal negation both
      belong to \(J\);
\item \emph{deductively closed relative to \(X\)} if
      \(J\models\varphi\) and \(\varphi\in X\) imply \(\varphi\in J\).
\end{enumerate}
Here \(\models\) is classical propositional semantic consequence. Write
\[
 \operatorname{cl}_X(B)=\{\varphi\in X:B\models\varphi\}
\]
for the agenda-relative consequences of any set \(B\) of formulas, and let
\(\mathcal J_X\) be the judgment sets satisfying all three requirements.
Thus \(J\) is closed exactly when \(\operatorname{cl}_X(J)=J\).

The source calls the second requirement simply ``consistency'' and explicitly
distinguishes it from satisfiability in its note 3
\cite[pp.~11, 27, 31]{ListPettit2001JudgmentPreprint}. Nonmembership
\(\varphi\notin J\) alone means nonassent. It implies acceptance of an opposing
proposition only when completeness supplies that additional commitment.
\index{judgment aggregation!completeness}\index{consistency!weak syntactic}
\index{deductive closure!relative to an agenda}
\end{hbdefinition}

\begin{hbtheorem}{thm:p17-ja-models}{proved=annex:p17.ja-models}{Rational judgment sets and truth assignments}
For any agenda with the opposing-proposition condition of
Definition~\ref{def:p17-ja-agenda},
\[
 J\in\mathcal J_X
 \quad\Longleftrightarrow\quad
 J=\{\varphi\in X:v\models\varphi\}
 \quad\text{for some classical valuation }v.
\]
Weak consistency together with agenda-relative closure already implies that
\(J\) is satisfiable. Weak consistency alone does not: the set
\(\{p,q,\neg(p\land q)\}\) contains no opposing pair but has no model.
The equivalence itself does not require the conjunction-containing subagenda.
\end{hbtheorem}

Closure here concerns only conclusions in \(X\). It is not a demand that a
finite collective record contain all formulas of the propositional language.
The theorem supplies coherent input sets by choosing valuations, including
values for atoms outside the displayed example.

\section{The doctrinal and discursive paradox}
\label{sec:p17-ja-paradox}

\begin{hbdefinition}{def:p17-ja-profile}{Profiles and proposition-wise majority}
A profile is an ordered tuple \(\mathbf J=(J_1,\ldots,J_n)\) of judgment sets.
For \(\varphi\in X\), define its support set and support count by
\[
 N_\varphi(\mathbf J)=\{i\in N:\varphi\in J_i\},
 \qquad s_\varphi(\mathbf J)=\#N_\varphi(\mathbf J).
\]
Strict proposition-wise majority returns
\[
 M(\mathbf J)=\{\varphi\in X:s_\varphi(\mathbf J)>n/2\}.
\]
This operation counts each proposition separately. Closure is an additional
operation, not part of this definition.
\index{majority rule!judgment aggregation}\index{discursive dilemma}
\end{hbdefinition}

\begin{hbexample}{def:p17-ja-doctrinal-profile}{Two premises and a rejected conjunction}
On \(X_0\), take the full truth judgments generated by these valuations:
\[
 \begin{array}{c|ccc}
 \text{individual}&p&q&p\land q\\ \hline
 1&1&1&1\\
 2&1&0&0\\
 3&0&1&0\\ \hline
 \text{strict majority}&1&1&0
 \end{array}
\]
Every omitted negation has the opposite individual truth value. The three
judgment sets are respectively
\[
 \{p,q,p\land q\},\qquad
 \{p,\neg q,\neg(p\land q)\},\qquad
 \{\neg p,q,\neg(p\land q)\}.
\]
This is the source's propositional profile
\cite[p.~12]{ListPettit2001JudgmentPreprint}. A doctrinal reading treats \(p\)
as the existence of a valid contract and \(q\) as the relevant breach, with
liability stipulated to require both. Voting on the two premises produces
liability; voting on the individual conclusions rejects it. The stipulated
doctrine is part of this example, not a claim about every legal rule.
\index{doctrinal paradox}\index{premise-based aggregation}
\index{conclusion-based aggregation}
\end{hbexample}

\begin{hbtheorem}{thm:p17-ja-paradox}{proved=annex:p17.ja-paradox}{Majority can destroy joint coherence}
The three sets in Example~\ref{def:p17-ja-doctrinal-profile} are in
\(\mathcal J_{X_0}\), but their majority set is
\[
 M(\mathbf J)=\{p,q,\neg(p\land q)\}.
\]
It is complete and weakly consistent, but neither satisfiable nor closed.
Its agenda-relative closure is all of \(X_0\), and is weakly inconsistent.

The disagreement can occur even when every individual rejects the conclusion.
For three premises \(p,q,r\), the valuations
\[
 (1,0,1),\quad(0,1,1),\quad(1,1,0)
\]
give majority support to each premise, while all three reject \(p\land q\land r\).
\end{hbtheorem}

In the source's workplace example, the premises concern a serious danger, an
effective safety measure and a bearable sacrifice. Each worker rejects a
different premise and therefore rejects the conjunction; the three separate
majorities endorse it \cite[pp.~7--8]{ListPettit2001JudgmentPreprint}. The issue
is the intersection of supporting coalitions. Large coalitions for separate
premises need not have a large intersection.

\section{Universal domain, anonymity and systematicity}
\label{sec:p17-ja-conditions}

\begin{hbdefinition}{def:p17-ja-uas}{Three conditions on an aggregation function}
A judgment aggregation function is a map
\(F:\mathcal D\to\mathcal P(X)\), where
\(\mathcal D\subseteq\mathcal J_X^n\) is its domain of rational profiles
and \(\mathcal P(X)\) is the set of subsets of \(X\).
The three source conditions are:
\begin{description}
\item[Universal domain (U).] \(\mathcal D=\mathcal J_X^n\): every rational
      profile is admissible.
\item[Anonymity (A).] For every admissible profile and every permutation
      \(\sigma:N\to N\),
      \[
       F(J_1,\ldots,J_n)=F(J_{\sigma(1)},\ldots,J_{\sigma(n)}).
      \]
\item[Systematicity (S).] One function \(f:\{0,1\}^n\to\{0,1\}\), fixed
      across propositions and profiles, satisfies
      \[
       \varphi\in F(\mathbf J)
       \quad\Longleftrightarrow\quad
       f(\delta_1(\varphi),\ldots,\delta_n(\varphi))=1,
      \]
      where \(\delta_i(\varphi)=1\) exactly when \(\varphi\in J_i\).
\end{description}
Systematicity combines dependence only on votes concerning the one proposition
with use of the same rule for every proposition. It is stronger than merely
requiring a function for each proposition separately. No monotonicity,
unanimity-preservation or majority threshold is assumed.
\index{judgment aggregation!universal domain}\index{anonymity}
\index{systematicity}
\end{hbdefinition}

\begin{hblemma}{thm:p17-ja-counts}{proved=annex:p17.ja-counts}{Anonymity makes the systematic rule depend on counts}
For the conjunction-containing agenda above, U, A and S give a function
\(g:\{0,\ldots,n\}\to\{0,1\}\) such
that \(f(d_1,\ldots,d_n)=g(\sum_i d_i)\). In particular, propositions with
equal support counts receive the same collective acceptance, even if their
supporting individuals differ.
\end{hblemma}

\begin{hbtheorem}{thm:p17-ja-impossibility}{proved=annex:p17.ja-impossibility}{List--Pettit impossibility for judgment aggregation}
Let \(n\geq2\), and let \(X\) satisfy the opposing-proposition condition and
contain distinct atoms \(p,q\) and \(p\land q\), as in
Definition~\ref{def:p17-ja-agenda}. There is no function
\[
 F:\mathcal J_X^n\longrightarrow\mathcal J_X
\]
satisfying U, A and S.

For even \(n\), the contradiction already comes from equal support for an
opposing pair, together with collective completeness and weak consistency.
For odd \(n\), equal support for \(p,q,\neg(p\land q)\) contradicts the
three collective rationality requirements. Thus ties are not the general cause
of the impossibility \cite[pp.~26--29, note~5]{ListPettit2001JudgmentPreprint}.
Permitting additional, nonrational input profiles cannot remove the
impossibility if all rational profiles remain admissible and A, S and the
same collective rationality requirements are retained.
\end{hbtheorem}

The proof in Annex~\ref{annex:p17.ja-impossibility} constructs the source's
profile for every \(n\), extends each row to a genuine valuation and handles
both possible outputs of the shared rule. It does not assume that the rule
follows a majority or even respects unanimous opinions. The conclusion concerns
the simultaneous requirements, rather than the impossibility of every coherent
group procedure.

\section{Restricting the domain: aligned judgments}
\label{sec:p17-ja-alignment}

\begin{hbdefinition}{def:p17-ja-aligned}{Unidimensional alignment}
A profile is \emph{unidimensionally aligned} if the individuals can be arranged
in a linear order such that, for every \(\varphi\in X\), the individuals
accepting \(\varphi\) form an initial segment or a final segment of that order.
Empty and full segments are allowed. The same order must work for all
propositions, although their accepting segments may point in different
directions. Different profiles may admit different orders.
\index{judgment aggregation!unidimensional alignment}
\end{hbdefinition}

\begin{hbtheorem}{thm:p17-ja-median}{proved=annex:p17.ja-median}{The median judgment set}
For odd \(n=2m+1\), strict majority on unidimensionally aligned profiles
\(\mathbf J\in\mathcal J_X^n\) returns
the entire judgment set of the middle individual in any witnessing order.
It therefore returns a member of \(\mathcal J_X\), and satisfies anonymity
and systematicity on this restricted domain.

For example, arrange five individuals as \(3,1,2,5,4\). Let the first three
accept \(p\) and reject \(q\), and let the final two reject \(p\) and accept
\(q\). All reject \(p\land q\). On \(X_0\), majority returns
\(\{p,\neg q,\neg(p\land q)\}\), the judgment set of individual 2.
\end{hbtheorem}

This is the source's alignment result and example
\cite[pp.~17--19]{ListPettit2001JudgmentPreprint}. Encouraging agreement is
not itself a proof that a profile has the required structure. The median
individual is selected by the profile's alignment; this is not the rule that
always appoints the same named person.

\section{Authority and the priority of premises or conclusions}
\label{sec:p17-ja-priority}

\begin{hbmethod}{def:p17-ja-priority-procedure}{Aggregate independent questions, then infer}
Choose a set \(Y\subseteq X\) with \(\neg y\in X\) for every \(y\in Y\).
Assume \(Y\) is logically independent in the following precise sense: every
assignment \(a:Y\to\{0,1\}\) extends to a classical valuation.
Let \(a(y)\) be the majority answer on \(y\), with a fixed anonymous tie rule
when needed, and form the signed set
\[
 B_Y(\mathbf J)=\{y\in Y:a(y)=1\}\cup
                \{\neg y:y\in Y,\ a(y)=0\}.
\]
Return \(F_Y(\mathbf J)=\operatorname{cl}_X(B_Y(\mathbf J))\).
Logical independence, rather than the label ``premise'', guarantees that the
chosen answers can be jointly satisfied.
\index{judgment aggregation!priority strategy}
\end{hbmethod}

\begin{hbtheorem}{thm:p17-ja-authority-priority}{proved=annex:p17.ja-authority-priority}{Which conditions authority and priority change}
On \(\mathcal J_X^n\), fixing an individual \(j\) and returning \(J_j\) gives rational collective
judgments and satisfies U and S, but fails A for \(n\geq2\).

The priority procedure of Definition~\ref{def:p17-ja-priority-procedure}
is defined on all rational profiles, is anonymous, and returns satisfiable,
weakly consistent, closed judgments. It is complete when every assignment to
\(Y\) determines the truth value of every formula in \(X\). Completeness need
not hold without that additional condition. The procedure can violate S.
\end{hbtheorem}

For Example~\ref{def:p17-ja-doctrinal-profile}, prioritize \(Y=\{p,q\}\).
The collective answers are \(p,q\), and their closure on \(X_0\) also accepts
\(p\land q\). Every issue is settled. If instead the group votes only on
\(c=p\land q\), it accepts \(\neg c\). On the conclusion-only agenda
\(\{c,\neg c\}\), this settles the whole question. On \(X_0\), its closure is
only \(\{\neg(p\land q)\}\): neither premise has been settled. Taking a
conclusion-driven decision therefore does not automatically supply a complete
judgment set on the larger agenda.

In a different case, a conclusion may have priority and constrain which premise
is rejected. The source describes this as choosing a modus-tollens orientation
instead of a modus-ponens orientation. There is no purely formal rule that
identifies the institutionally appropriate priorities in every setting
\cite[pp.~20--23]{ListPettit2001JudgmentPreprint}. An authority called in only
when a paradox occurs may alter systematicity as well as anonymity.

\section{Special support and incomplete collective judgments}
\label{sec:p17-ja-special-support}

\begin{hbtheorem}{thm:p17-ja-unanimity}{proved=annex:p17.ja-unanimity}{Unanimity preserves consistency and closure}
For rational inputs, the unanimity rule
\[
 U_\cap(\mathbf J)=\bigcap_{i\in N}J_i
\]
is satisfiable, weakly consistent and closed. It satisfies U, A and S, with
\(f(d)=1\) exactly when every coordinate of \(d\) is 1. It need not be complete.
Taking its agenda-relative deductive closure adds nothing.
\end{hbtheorem}

\begin{hbdefinition}{def:p17-ja-finite-issues}{A finite agenda counted by issue pairs}
For the following support bound, specialize to a finite agenda
\(X=\{\theta_1,\neg\theta_1,\ldots,\theta_k,\neg\theta_k\}\), where the
displayed formulas are distinct and each formula belongs to exactly one
displayed pair. Thus \(k\) counts opposing-proposition pairs, not \(2k\)
separate strings. In the running agenda \(X_0\), \(k=3\).
Define the strict supermajority rule
\[
 H_k(\mathbf J)=
 \{\varphi\in X:s_\varphi(\mathbf J)>(k-1)n/k\}.
\]
The strict inequality implements the source's requirement of support
\emph{more than} \((k-1)/k\) \cite[pp.~23--24]{ListPettit2001JudgmentPreprint}.
\index{judgment aggregation!supermajority}
\end{hbdefinition}

\begin{hbtheorem}{thm:p17-ja-supermajority}{proved=annex:p17.ja-supermajority}{A sufficient support bound and the effect of closing it}
Assume the finite agenda above, \(k\geq2\), and rational individual judgments.
Then \(H_k\) satisfies U, A and S, and its accepted propositions have a common
model. It need not be complete or deductively closed. The rule
\[
 G_k(\mathbf J)=\operatorname{cl}_X(H_k(\mathbf J))
\]
is satisfiable, weakly consistent and closed, and satisfies U and A. It need not
be complete or systematic.

The strict threshold is essential to this general guarantee. On \(X_0\) with
three individuals, allowing support \emph{at least} \(2/3\) admits precisely
the jointly unsatisfiable majority set in Theorem~\ref{thm:p17-ja-paradox}.
\end{hbtheorem}

Unanimity can leave a group silent on disputed issues. A lower special-support
threshold can permit more conclusions, but taking further consequences can make
an issue's output depend on votes concerning other issues. The support bound
and the closure operation therefore discharge different obligations.

\section{The six escape routes and their boundaries}
\label{sec:p17-ja-escapes}

The source distinguishes six departures from the incompatible package
\cite[pp.~16--25]{ListPettit2001JudgmentPreprint}:
\begin{enumerate}
\item \emph{Convergence:} restrict U to profiles with an appropriate common
      structure; odd-electorate alignment is a proved sufficient condition.
\item \emph{Authority:} relax A by assigning a named individual's judgments
      special authority. A conditional authority may also change S.
\item \emph{Priority:} relax S by settling selected independent questions first
      and inferring answers elsewhere. Completeness also needs a determining
      premise set or an additional completion procedure.
\item \emph{Special support:} relax collective completeness. Unanimity retains
      closure and S; strict supermajorities retain satisfiability, while their
      closure may sacrifice S.
\item \emph{Relax consistency:} retain enough accepted propositions and
      consequences to include opposing judgments. Closing an inconsistent
      majority set is one such procedure, and may also sacrifice S.
\item \emph{Relax closure:} retain separate majority judgments without their
      consequences. With odd electorates this preserves completeness and weak
      consistency, but not necessarily satisfiability.
\end{enumerate}

\begin{hbtheorem}{thm:p17-ja-boundaries}{proved=annex:p17.ja-boundaries}{Explicit boundaries of the incompatible package}
For odd electorates, raw majority satisfies U, A and S and produces complete,
weakly consistent judgments, but can fail closure. Closing the source's
paradoxical majority set makes it inconsistent. The constant rule
\(F(\mathbf J)=X\) is complete, closed and satisfies U, A and S, but is
inconsistent on the present agendas.

With one individual the identity rule satisfies the whole rationality package.
On an agenda consisting only of independent atoms and their negations,
odd-electorate majority also returns a rational set. Neither case meets all
the impossibility theorem's hypotheses.
\end{hbtheorem}

The final two strategies show why the source regards weak syntactic coherence
as insufficient to guide collective action. Conversely, an incomplete but
satisfiable record can represent genuine suspension of judgment. Calling both
outcomes ``no collective decision'' would conceal their logical difference.

\begin{hbdefinition}{def:p17-ja-relative-systematicity}{Systematicity within proposition classes}
For an equivalence relation \(\sim\) on \(X\), a weaker condition permits a
separate function \(f_C:\{0,1\}^n\to\{0,1\}\) for each equivalence class
\(C\), fixed across profiles. Acceptance of \(\varphi\in C\) is then computed
by \(f_C\) from the votes on \(\varphi\). Ordinary S is the case with one
class. The source raises structure-dependent versions in note 4; the theorem
proved here concerns ordinary S and does not assert the same impossibility for
every choice of \(\sim\).
\end{hbdefinition}

Committees, courts and associations may combine these strategies. The same
problem also arises over time: a new collective answer may conflict with
commitments already adopted. A procedure must state which recorded commitments
are retained, which may be revised, whose judgments have authority and which
questions remain undecided. The impossibility theorem identifies a conflict
among logical and procedural requirements; it does not select the appropriate
institutional compromise or guarantee that deliberation produces convergence.
